\section{Summary}\label{sec:1}

This system runs demand-controlled ventilation over six rooms of an office floor and, on top of
it, a second loop that decides whether that ventilation is still \emph{working}. The guiding
principle is one sentence and it is the thread through every decision that follows: \textbf{a
stuck damper and a full room look identical to a CO₂ controller, and what tells them apart is a
model of what the last command should have achieved.} A controller that sees only concentration
cannot distinguish a room that will not ventilate from a room with more people in it; both read
high. Predicting the response each command owes, and watching the gap between that prediction and
the measurement, turns the fault from an invisible condition into an observable one, which makes
the autonomous part an inference problem rather than a threshold, and makes fault injection and
degraded operation the use case itself rather than an experiment bolted on at the end.

The system is therefore \textbf{two nested loops}, and the outer one is the contribution.

\begin{longtable}[]{@{}
  >{\raggedright\arraybackslash}p{(\columnwidth - 4\tabcolsep) * \real{0.26}}
  >{\raggedright\arraybackslash}p{(\columnwidth - 4\tabcolsep) * \real{0.10}}
  >{\raggedright\arraybackslash}p{(\columnwidth - 4\tabcolsep) * \real{0.64}}@{}}
\toprule
\textbf{Loop} & \textbf{Period} & \textbf{Closes because} \\
\midrule
\endhead
Inner, control (\texttt{dcv-controller}) & ≈ 1 s & A damper command changes airflow, which
changes CO₂, which the controller reads next \\
Outer, supervision (\texttt{fault-monitor} + \texttt{supervisor}) & ≈ 60 s & A confirmed fault
changes the control strategy, which changes the commands, which changes the residual the monitor
reads next \\
\bottomrule
\end{longtable}

The inner loop is deliberately a plain rule: CO₂ in, damper position out, with hysteresis and a
freshness check. That is not an omission but the point. Demand-controlled ventilation is a solved
problem, and making the inner rule cleverer would have bought nothing the course grades; the
interesting and unsolved question is what happens when the actuator the clever rule depends on
has quietly stopped working. Keeping the inner loop simple also keeps it \emph{verifiable}, which
matters because the outer loop judges it: an observer that has to model a sophisticated
controller inherits that controller's complexity as model error.

\begin{longtable}[]{@{}
  >{\raggedright\arraybackslash}p{(\columnwidth - 2\tabcolsep) * \real{0.26}}
  >{\raggedright\arraybackslash}p{(\columnwidth - 2\tabcolsep) * \real{0.74}}@{}}
\toprule
\textbf{Field} & \textbf{Value} \\
\midrule
\endhead
\textbf{Project title} & Ventilation Fault Detection, Diagnosis and Graceful Degradation \\
\textbf{Team members} & Wassim Akachi (project group 5, solo) \\
\textbf{Use case} & Demand-controlled ventilation that detects, diagnoses and survives faults in
its own equipment \\
\textbf{Repository} & \url{https://github.com/eduni-pro/202608-d7065e}, subfolder \texttt{project/} \\
\textbf{Version / date} & v1.1, submission, 2026-10-11 \\
\bottomrule
\end{longtable}

This is a solo project, so no contribution statement is included: the course requires one only
where a pair submits a single system, and every decision, measurement and line of code recorded
here is my own.

\textbf{Where it landed, before the detail.} Two of the seven requirements are met outright. Three
more pass on the arm their acceptance criterion names and fail on the harder arm beside it, and
they are reported as \emph{two} numbers rather than one, because a single verdict would have to
pick which regime to be silent about. \textbf{R-3 is the one that fails everywhere (0.55
confirmed false alarms per fault-free simulated week against a budget of one, and a worst week of
6, which is the criterion it is scored on) and R-1's, R-2's and R-7's harder arms all have that
same false-alarm rate somewhere in their causes.} One defect,
surfacing four times, is a more useful thing to know about this system than four defects would be,
and §\ref{sec:11} reports it that way rather than softening it. The regime that separates the arms is not
the same in each case and is named with each result: for R-2 and R-7 it is \emph{occupancy} (the
drawn weekday passes, design occupancy does not), while for R-1 it is the \emph{protocol}, which
detects 20 of 20 at a p95 of 6 min when the damper is stepped deliberately and 15 of 20 at 257 min
when the fault must be found in the closed loop that is also compensating for it.

\subsection{Changes since the proposal}\label{sec:1.1}

Three things changed, and each was forced by a decision I had to make rather than by a change of
mind. The accepted proposal (2026-09-18) carried one examiner condition, \emph{``make sure to
include several rooms''}, and answering it honestly is what drove all three.

\textbf{1. The coupling topology: six independent rooms became one shared air handler with a
damper per room.} ``Several rooms'' can be satisfied by replication, six copies of the single-room
spike with nothing one room can do to another, and the implementation plan's risk register names
exactly that as the failure to avoid. ADR-0008 therefore weighed three couplings and chose
\textbf{option B}: one air-handling unit \texttt{AHU-A120}, one damper per room, and a
proportional flow split, sized 43 \% above the sum of the six rooms' design flows, but only
0.4 \% above it at the normal fan speed of 0.7, so that \emph{the only way one room gets more air
is if another damper closes} and the headroom is a compensation budget rather than normal slack. Rejected: option A
(independent rooms, a fan each), which makes R-7 hold by construction and its experiment measure
nothing; and option C (option B plus inter-room air exchange), which is the proposal's literal
wording but whose exchange coefficient, the one parameter nothing can measure, would be what
decides whether R-2 passes.

\textbf{2. ``Over-ventilate the neighbour'' became ``over-drive the shared handler and trim the
neighbours' dampers''.} The proposal said the supervisor would compensate by over-ventilating the
faulted room's neighbour, which presumes air moves between rooms. Under option B it does not:
rooms are coupled through the \emph{duct}, not through the wall. The mechanism that replaces it is
better, not merely different, and Section \ref{sec:7.4} derives it: one common trim on the healthy dampers
plus one raise of the shared fan speed holds every healthy room's airflow exactly where it was
while multiplying the faulted room's flow by the reciprocal of the trim. The cost is a stated
bound (a damper stuck fully shut has no air path and cannot be reached at all), and that bound is
reported with R-2 rather than hidden.

\textbf{3. Fan power moved from a per-room signal to a plant signal, and that is what makes
discrimination possible at all.} In the proposal's sketch each room had its own fan, so fan power
was a local quantity and a blocked filter was a local anomaly. With one shared unit, fan power
belongs to the plant. That single relocation is what gives the monitor a \emph{cross-room} signal:
a blocked filter starves all six rooms at once and raises power per unit of flow, while a stuck
damper moves one room's residual and leaves plant power flat. Without it, separating those two
faults would have needed a fourth invented signal, which the risk register explicitly forbade.

One thing did \emph{not} change and deserves saying so: the detection method. The proposal
proposed a predicted step response, a residual and a CUSUM (one-sided then, two-sided now that
§\ref{sec:7.3}'s guarded lower arm has been added), and the week-37 spike
(Section \ref{sec:11.1}) found the separation so large that the design question moved from ``can this be
detected'' to ``how few false alarms can it be detected with''.\label{sec:1end}
